%% ============================================================
%% 分片 15：原文第 6 节「Loop and Line Operators」（原文 glshort.tex 起于第 4483 行）
%% 由 tools/harvmac2tex.py 从 arXiv 源码（hep-th/0604151v3）自动生成。
%% 只译英文散文；公式、编号、\label、\cite 一律不改。
%% ============================================================

\section{Loop and Line Operators}

\label{sec:loopop} 

 What is really unusual about the four-dimensional TQFT's
introduced in section 第~\ref{sec:twisting}~节, compared to other theories with a
superficially similar origin, is that they admit operators that
are associated to oriented one-manifolds ${\cal S}\subset M$, and
depend on the homotopy class, not just the homology class, of
$\cal S$. The point is not that we want to study the fundamental
group of $M$, but that these one-manifold operators prove to have
very interesting properties in relation to branes.

If $M$ is compact and without boundary, then $\cal S$ is required
to be a closed loop. Reflecting this case, one-manifold operators
are commonly called loop operators.  More generally, however, on a
noncompact  four-manifold, we allow line operators that go off to
infinity. An important example is a static line operator, whose
trajectory spans all time at a given point in space.  Such a line
operator must be included in quantizing the theory to construct a
Hilbert space of physical states.  This will actually be the most
important case in our applications.\footnote{Thus, other cases such as
line operators that end on branes, though they make sense in
general, will not be important in the present paper.}

 
 

In section 第~\ref{sec:topwil}~节, we discuss the most elementary loop or line
operators, usually called Wilson  operators. Wilson operators in
TQFT's similar to the ones we consider here, with the connection
modified by a scalar field to ensure the topological symmetry,
were first constructed by Blau and Thompson \cite{ref84}. Similar
gauge-invariant operators were also used in studying
five-dimensional supersymmetric Yang-Mills theory \cite{ref85}. Such
operators are important in contemporary developments in string
theory (along with their magnetic duals, which we introduce in
section 第~\ref{sec:tophooft}~节). They were first introduced in the half-BPS
case in \cite{ref86,ref87}. The topological line operators that we
consider are most similar to the 1/16 BPS line operators defined
in \cite{ref88} for any
loop ${\cal S}\subset \Bbb{R}^4$.

In section 第~\ref{sec:tophooft}~节, we describe the dual 't Hooft  operators. In
section 第~\ref{sec:comptwo}~节, we discuss the interpretation of such operators
after compactification to two dimensions. Finally, in section
第~\ref{sec:lineop}~节, we study line operators as operators on branes.

\subsection{Topological Wilson Operators}

\label{sec:topwil}

In gauge theory, the most elementary loop operator is the Wilson
loop operator.  We consider gauge theory with gauge group $G$ and
connection $A$ on a $G$-bundle $E\to M$. We let ${\cal S}$ be an
oriented loop in $M$, oriented say by the choice of a one-form
$ds$, where $s$ is a parameter along ${\cal S}$.  We pick an
irreducible representation $R$ of $G$, and let\footnote{Since our
covariant derivative is $D=d+A$, the holonomy operator has a minus
sign in the exponent.  In terms of a hermitian gauge field $A'=iA$
(recall eqn. \eqref{eq:boloko}), the holonomy operator is ${\rm
Tr}_R\,P\,\exp\left(i\oint_{\cal S} A'\right)$.}
% [WARN] 公式 \eqn\eqref{eq:ytro} 无编号或重复
\begin{equation*}
\tag*{(6.1)}
W_0(R,{{\cal S}})={\rm Tr}_R
\,P\,\exp\left(-\oint_{{\cal S}} A\right).
\end{equation*}
 In other words,
$W_0(R,{\cal S})$ is the trace, in the representation $R$, of the
holonomy of $A$ around ${\cal S}$.

Now let us look back to eqn. \eqref{eq:bono}\ from section 第~\ref{sec:twisting}~节, in
which we determined the transformation law of $A$ under the
topological symmetry: % 式 (6.2)
\setcounter{equation}{1}
\begin{equation}\label{eq:likp}
\delta_T
A_\mu=iu\psi_\mu+iv\tilde\psi_\mu.
\end{equation}
 Clearly, for any non-zero
$(u,v)$, the Wilson loop operator is not invariant.  However,
twisted ${\EUN}=4$ Yang-Mills theory also has an adjoint-valued
one-form $\phi=\sum_{\mu=0}^3\phi_\mu\,dx^\mu$, transforming as
% 式 (6.3)
\setcounter{equation}{2}
\begin{equation}\label{eq:nikp}
\delta_T\phi_\mu=iv\psi_\mu-iu\tilde\psi_\mu.
\end{equation}
 It is
possible to modify the connection $A$ by adding to it a multiple
of $\phi$.

Precisely if $t=v/u$ is equal to $\pm i$, a linear combination of
$A$ and $\phi$ is invariant.  For $t=i$, the invariant combination
is $\CA=A+i\phi$; at $t=-i$, it is $\overline\CA=A-i\phi$.  We
then define the supersymmetric Wilson loop operator as the
holonomy of $\CA$ or $\overline\CA$ around ${\cal S}$.  Thus
% 式 (6.4)
\setcounter{equation}{3}
\begin{equation}\label{eq:ymop}
W(R,{\cal S})=\Tr_R\,P\exp\left(-\oint_{{\cal
S}}\CA\right)=\Tr_R\,P\exp\left(-\oint_{{\cal
S}}(A+i\phi)\right),~~ t=i
\end{equation}
 and % 式 (6.5)
\setcounter{equation}{4}
\begin{equation}\label{eq:jymop}
W(R,{{\cal S}})=
\Tr_R\,P\exp\left(-\oint_{{\cal S}}{\overline\CA}\right)
=\Tr_R\,P\exp\left(-\oint_{{\cal S}}(A-i\phi)\right), ~~t=-i.
\end{equation}
 The
central generator $\EUI$ of $SU(4)_{{\cal R}}$, which acts by
$t\to -t$ and $\phi\to -\phi$, transforms one formula into the
other.

The existence of Wilson loop operators possessing the topological
symmetry is extremely natural at $t=\pm i$.  At these values of
$t$, the topological equations \eqref{eq:loopme}\ assert the flatness of
$\CA$. So the holonomies  of $\CA$ or $\bar\CA$ around closed
loops are natural invariants. Equally well, such  holonomies are
not natural invariants if $t\not=\pm i$, since the topological
equations do not necessarily imply flatness.  (The vanishing
theorems of section 第~\ref{sec:vanishing}~节, which relate supersymmetric
configurations to complex flat connections for generic $t$,
require some global input and do not hold in all situations, for
instance involving branes, to which one might apply the TQFT.)

Though we formulated these definitions for ${{\cal S}}$ a closed
loop, they have a good analog for the open case. If ${{\cal S}}$
has endpoints $p$ and $q$, we define $W(R,{{\cal S}})$ not  as a
trace but as the matrix of parallel transport (of the connection
$\CA$ or $\bar\CA$) from the fiber of $E$ at $p$, taken in the
representation $R$, to the fiber at $q$. $W(R,{{\cal S}})$ is thus
a map from a vector space at $p$ to a vector space at $q$.  It is
then incorporated as part of a larger quantum construction
involving initial and final quantum states, as mentioned in our
introductory remarks. This extended version of the Wilson operator
will be essential when we get to the geometric Langlands program
in earnest.

  When the
context is clear, we abbreviate $W(R,{{\cal S}})$ as $W(R)$ or
simply $W$.
